% =====================================================================
%  Section 12 — Phase 2 research questions Q13–Q25
% =====================================================================
\makesectioncover{Phase 2 \textbullet\ The Thirteen Questions}

\section{Phase 2 \textbullet\ Thirteen Research Questions}
\label{sec:phase2-questions}

The Phase 2 analysis notebook poses thirteen questions, numbered Q13
to Q25 to continue the Phase 1 numbering. Each is answered with a
single-method or multi-method analysis on the integrated dataset
plus the cleaned Phase 1 outputs. We use the same five-part template
as Section~\ref{sec:phase1-questions}: \emph{Why}, \emph{Data},
\emph{Method}, \emph{Visualisation \& reading}, \emph{Why this viz}.

% --------------------------------------------------------------------
\subsection{Q13 \textbullet\ Metro coverage \texttimes\ density}

\textbf{Why.} The first descriptive question of Phase 2: did metro
expansion follow density? Newer lines should plausibly serve denser
districts than older ones if the network was demand-aligned.

\textbf{Data.} \fn{metro\_stations.csv} (S3, 89 stations with opening
year and line) joined to \fn{districts\_wide.csv} (S6, 68 districts with
population) at the closest district for each station.

\textbf{Method.} Group stations by line; for each line plot opening
year against the population of the closest district. Spearman
correlation \emph{within} each line lets the reader see whether the
relationship has different signs for the older lines (1, 2) vs the
newer ones (3, LRT).

\figitem{phase2/q13_coverage_vs_density.png}{1.0}{Q13 \textbullet\ Per-line scatter of station opening year vs. closest-district population.}{fig:q13a}

\textbf{Why this viz.} A faceted scatter is the right form when the
question is multi-line. The line-by-line breakdown isolates the
diverging trends; a single pooled scatter would hide them.

\textbf{Reading.} Line 1 (the oldest) sits in the densest districts;
Line 3 and the LRT have their newer stations in lower-density
districts. The visual signature of the Phase 2 story.

% --------------------------------------------------------------------
\subsection{Q14 \textbullet\ Ghost terminals near the new metro}

\textbf{Why.} The most direct Phase 1 $\to$ Phase 2 bridge. Phase 1
identified 115 ghost terminals; Q14 asks how many of them are within
walking distance of a post-2014 metro station. The expectation, in a
demand-aligned network, would be high.

\textbf{Data.} \fn{merged\_G\_underused\_terminals.csv} (Phase 1, 115
ghosts) and the post-2014 cohort of \fn{metro\_stations.csv} (Phase 2).

\textbf{Method.} For each ghost, haversine distance to every post-2014
metro station; keep the minimum; bucket
($0$--$1$\,km, $1$--$2$, $2$--$5$, $5$--$10$, $>$$10$\,km).

\figitem{phase2/q14_distance_buckets.png}{1.0}{Q14 \textbullet\ Distance from each Phase 1 ghost terminal to the nearest post-2014 metro station.}{fig:q14a}
\figitem{phase2/q14_spatial_diagnostic.png}{1.0}{Q14 \textbullet\ Spatial diagnostic — ghosts and new metro on the same map.}{fig:q14b}

\textbf{Why this viz.} The bucketed bar chart is the right form when
the message is categorical (\emph{how many ghosts in each distance
band}). The companion map gives the same information geographically
for readers who think in shapes rather than histograms.

\textbf{Reading.} \stat{Only 9 of 115 ghosts (8\%)} sit within 1\,km of
any post-2014 metro station. \stat{97 of 115 (84\%)} are beyond 2\,km.
The buckets read \stat{3 / 6 / 9 / 27 / 70} from nearest to farthest.
The new metro did not relieve old underused terminals.

% --------------------------------------------------------------------
\subsection{Q15 \textbullet\ Metro $\to$ terminal integration over time}

\textbf{Why.} Even where new metro and old terminals do coexist, do they
\emph{integrate}? Q15 asks whether each new metro line opened near
existing Phase 1 terminals (so transfer is short) or in greenfield
territory (so it functions as its own ecosystem).

\textbf{Data.} Metro stations (S3, with opening year) + terminals
(Phase 1).

\textbf{Method.} For each metro station compute the haversine to its
nearest terminal; group by line; plot distributions and time-series.

\figitem{phase2/q15_metro_terminal_integration.png}{1.0}{Q15 \textbullet\ Per-line transfer distance distribution. Line~1 stations are typically within 200\,m of an existing mawfaq; Line 3's newer stations are not.}{fig:q15}

\textbf{Why this viz.} A distribution per line is more informative than
a mean — the question is about the shape (does the line consistently
sit near terminals or are there long-tail outliers).

\textbf{Reading.} Older lines integrate tightly with the informal
network; the post-2012 expansions sit further from existing
terminals. The integration story degraded over time.

% --------------------------------------------------------------------
\subsection{Q16 \textbullet\ Fastest-growing districts \& coverage}

\textbf{Why.} A planning network that follows demand should serve
\emph{growth} as much as standing population. Q16 isolates the 15
districts with the highest CAGR (2006$\to$2017) and asks whether they
have any new-mode coverage at all.

\textbf{Data.} \fn{districts\_wide.csv} (S6) + integrated station layer.

\textbf{Method.} CAGR rank top-15; per-district station count within a
3-km radius; lollipop chart (CAGR on x, stations on y).

\figitem{phase2/q16_cagr_slope.png}{1.0}{Q16 \textbullet\ Top-15 fastest-growing districts vs. number of new-mode stations within 3\,km.}{fig:q16}

\textbf{Why this viz.} A lollipop chart is the right form when one of
the two axes is highly compressed (most districts have 0--3 stations).

\textbf{Reading.} Several of the highest-CAGR districts have zero
new-mode coverage. Growth is not driving infrastructure placement.

% --------------------------------------------------------------------
\subsection{Q17 \textbullet\ Density vs underserved score}

\textbf{Why.} Tests whether ``underservedness'' (Phase 1 G4) is simply
a density artefact or a real coverage mismatch. If it were just
density, every dense hex would be underserved and the score would
correlate trivially with population. If the score \emph{also}
correlates with population, that is the signal that the underservedness
is structural — these are the districts where the gap is largest.

\textbf{Data.} \fn{q17\_dense\_underserved\_hexes.csv} (recomputed from
Phase 1's hex layer + S6 district population).

\textbf{Method.} Spearman correlation; joint density / hex-bin scatter
with marginal histograms; quadrant overlay highlighting ``dense AND
underserved''.

\figitem{phase2/q17_density_underserved.png}{1.0}{Q17 \textbullet\ Population density vs underserved score. Spearman $\rho = 0.548$, $p < 10^{-100}$.}{fig:q17}
\figitem{phase2/q17_target_tier_heatmap.png}{1.0}{Q17 \textbullet\ Decile heat-map of dense-and-underserved hexes. The top-right is the policy target.}{fig:q17b}

\textbf{Why this viz.} A joint density / hex-bin plot with marginals is
the canonical form for a ``where do points pile up?'' question. The
quadrant overlay does the storytelling work in-figure.

\textbf{Reading.} \stat{$\rho = 0.548$}, \stat{$p < 10^{-100}$}. Densest
hexes are also the most underserved. The market problem is not just
density; it is density \emph{plus} weak coverage. This is the
correlation that the K-Means segmentation in Q24 will use as a feature.

% --------------------------------------------------------------------
\subsection{Q18 \textbullet\ Informal share by density}

\textbf{Why.} If informal microbus dependence concentrated only in the
dense core, it would be a poverty-pocket problem, not a structural
one. Q18 tests whether the informal share is correlated with density
across the city.

\textbf{Data.} Phase 1 boarding (informal vs total) joined to S6
district population.

\textbf{Method.} Spearman + OLS slope.

\figitem{phase2/q18_informal_share.png}{1.0}{Q18 \textbullet\ Informal-mode share by hex population. The trendline is essentially flat.}{fig:q18a}
\figitem{phase2/q18_per_tier_box.png}{1.0}{Q18 \textbullet\ Informal share distribution by density tier.}{fig:q18b-tier}

\textbf{Why this viz.} A scatter with regression line shows both the
spread and the flatness; the per-tier box plot rules out the
possibility that the flatness is an aggregation artefact.

\textbf{Reading.} \stat{$\rho = 0.025$}, \stat{$p = 0.40$}.
Statistically flat. Informal transport averages \stat{$\sim$47\%} mode
share \emph{across all density tiers}. Microbuses are not a
poverty-pocket phenomenon — they are the everyday transport substrate.
This single result is the structural argument behind every later
``two systems'' framing.

% --------------------------------------------------------------------
\subsection{Q18b \textbullet\ Gap-closure matrix · the verdict}

\textbf{Why.} The single most important question of Phase 2. We
cross-tabulate the four Phase 1 gaps (G1--G4) against the three
post-2014 modes (Metro L3, LRT, BRT). Each cell is the percentage of
that gap's sites within 2\,km of a station of that mode.

\textbf{Data.} The four gap-site sets (115 ghosts, 19 empty-return,
$\sim$1{,}340 vehicle-mismatch routes, 79 underserved hexes) buffered
against each mode's stations.

\textbf{Method.} For each (gap, mode) cell: count gap-category sites
within a 2-km haversine buffer of any station of that mode; report as
a percentage.

\figitem{phase2/q18b_matrix.png}{1.0}{Q18b \textbullet\ The gap-closure
matrix. \stat{No cell above 25\%}; most below 16\%; the single best
cell is BRT $\times$ G3 (vehicle mismatch) at \stat{25.0\%}. This is
the empirical verdict on the \$10B build-out.}{fig:q18b}

\textbf{Why this viz.} A heatmap is the only form that lets the eye
scan all twelve cells at once. The annotated values mean we don't need
a separate table.

\textbf{Reading.} The $3 \times 4$ matrix shows that:
\begin{itemize}
  \item Metro L3 partially relieves G1 ghosts (15.7\%) but barely
        touches the other three gaps.
  \item LRT closes G3 (vehicle mismatch) at 25\% — its only above-token
        contribution.
  \item BRT is the closest the new infrastructure comes to a planning
        win: 25\% on G3, 13.1\% on G4 (underserved hexes).
\end{itemize}
\textbf{The single sentence verdict: \stat{the post-2014 infrastructure
closes a quarter or less of any single gap, and most cells are below
sixteen percent}.}

% --------------------------------------------------------------------
\subsection{Q19 \textbullet\ GTFS coverage — formal vs informal data}

\textbf{Why.} Beyond the physical-infrastructure gap there is a
\emph{data} gap: how much of the network is published in GTFS at all?
If the formal system has full GTFS coverage and the informal system
has none, that is itself a product opportunity (Section 16).

\textbf{Data.} \fn{gtfs\_stops.csv} (S1) vs Phase 1 boarding.

\textbf{Method.} For each Phase 1 stop, find the closest GTFS stop;
flag covered if within 800\,m; aggregate by agency.

\figitem{phase2/q19_agency_pie.png}{0.85}{Q19 \textbullet\ GTFS agency composition.}{fig:q19a}
\figitem{phase2/q19_gtfs_coverage.png}{1.0}{Q19 \textbullet\ Phase 1 stops covered vs uncovered by published GTFS.}{fig:q19b}

\textbf{Why this viz.} A pie shows agency composition cleanly; the
coverage map answers the spatial question.

\textbf{Reading.} A large fraction of informal-heavy stops sit beyond
the 800-m GTFS catchment. The data asymmetry is real and quantifiable.

% --------------------------------------------------------------------
\subsection{Q20 \textbullet\ BRT corridor vs informal demand}

\textbf{Why.} The BRT is the H3 contrast case — Q20 is its descriptive
warm-up. We rank every BRT station by the informal demand inside its
500-m buffer and ask whether the ranking is dominated by stations sited
on \emph{existing} demand.

\textbf{Data.} \fn{brt\_stations.csv} (S5) + Phase 1 boarding.

\textbf{Method.} For each BRT station, sum daily informal boardings
within a 500-m buffer; rank descending; horizontal bar.

\figitem{phase2/q20_brt_corridor.png}{1.0}{Q20 \textbullet\ BRT stations ranked by nearby informal demand. The top-right tail is the corridor that did inherit existing demand.}{fig:q20}

\textbf{Reading.} A clear top-tier of BRT stations carries hundreds of
informal boardings within 500\,m. The bottom tail does not — but the
top tier is what drives the H3 result.

% --------------------------------------------------------------------
\subsection{Q21 \textbullet\ Fare per kilometre — formal vs informal}

\textbf{Why.} Accessibility is not only spatial; it is economic. Q21
asks how much riders \emph{pay per km} on each mode.

\textbf{Data.} \fn{cleaned\_routes.csv} (Phase 1) + GTFS fare attributes
(S1).

\textbf{Method.} For each route, $\text{LE/km} = \text{Fare} /
\text{Route\_Length\_km}$; group by mode/vehicle type; box plot.

\figitem{phase2/q21_fare_per_km.png}{1.0}{Q21 \textbullet\ Fare per
kilometre by mode. Microbus and tomnaya carry the informal premium.}{fig:q21}

\textbf{Reading.} Formal bus = \stat{0.22 LE/km}, formal metro =
\stat{0.25 LE/km}. Informal microbus = \stat{0.62 LE/km}
($\sim$3$\times$ formal). Informal tomnaya = \stat{1.30 LE/km}
($\sim$6$\times$). Riders pay the premium because the formal alternative
does not reach them.

% --------------------------------------------------------------------
\subsection{Q22 \textbullet\ Metro under-performance relative to population}

\textbf{Why.} Move from raw counts to coverage \emph{expectation}. Fit
metro stations as a function of population per district; the
\emph{residual} (actual minus predicted) is the under-coverage
metric.

\textbf{Data.} \fn{metro\_stations.csv} (S3) + S6 districts.

\textbf{Method.} OLS regression of stations on population (3-km radius
per district centroid); rank residuals.

\figitem{phase2/q22_residual_ranked.png}{1.0}{Q22 \textbullet\ Top-15 most under-served districts by metro residual.}{fig:q22a}
\figitem{phase2/q22_governorate_box.png}{1.0}{Q22 \textbullet\ Residual distribution by governorate. Giza/Qalyubia carry the worst tail.}{fig:q22b}

\textbf{Reading.} The four worst residuals are
\strval{Ad-Duqqi} ($-9.28$), \strval{Sheikh Zayed} ($-9.09$),
\strval{6th October City 1 \& 3} ($-8.52$),
\strval{Al-Hawamidiyah} ($-8.14$). All in Giza/Qalyubia — dense or
rapidly-growing edges. \stat{This is the ranking that powers the
headline coverage-need map.}

% --------------------------------------------------------------------
\subsection{Q23 \textbullet\ Adly Mansour as a convergence node}

\textbf{Why.} Adly Mansour is symbolically important — the Transport
Minister has framed it as Cairo's ``seventh mode'' interchange. Q23
quantifies whether the symbolic claim matches the demographics.

\textbf{Data.} Integrated layer (metro, LRT, BRT, Phase 1 informal) +
S6 district population, in a 2.5-km cluster centred on Adly Mansour.
Compared against 150 random 2.5-km clusters elsewhere in Cairo.

\textbf{Method.} For each random cluster compute station density
(stations + terminals within 2.5\,km); rank Adly Mansour against the
empirical distribution.

\figitem{phase2/q23_modal_pie.png}{0.85}{Q23 \textbullet\ Modal composition within 2.5\,km of Adly Mansour.}{fig:q23a}
\figitem{phase2/q23_percentile_histogram.png}{1.0}{Q23 \textbullet\ Adly Mansour vs 150 random Cairo 2.5-km clusters by station density.}{fig:q23b}
\figitem{hero/adly_mansour_zoom.png}{1.0}{Adly Mansour zoom — case-study map of the multimodal node.}{fig:adly-zoom}

\textbf{Reading.} \stat{4 modes within 2.5\,km}, \stat{8
stations/terminals total}, but only the \stat{21st percentile} of
random Cairo clusters by density. Symbolically multimodal,
demographically average.

% --------------------------------------------------------------------
\subsection{Q24 \textbullet\ K-Means coverage-need synthesis (AI Technique 2)}

\textbf{Why.} Single-variable rankings cannot describe a market. The
project needs \emph{segmentation} — groups of districts that share a
joint profile of demand, density, growth, and infrastructure. Q24
delivers it via K-Means on five engineered features.

\textbf{Data.} S6 districts + integrated station layer.

\textbf{Method.} K-Means with \stat{$k=4$}, k-means++ init, $n_{\text{init}}=50$,
features = \fn{log\_pop\_2023}, \fn{CAGR \%}, stations within 3\,km,
informal stops within 3\,km, Cairo-governorate indicator.

\textbf{Validation.} \stat{Adjusted Rand Index = 1.00 across 5 random
seeds} — the partition is perfectly stable. Silhouette at $k=4$ is
\stat{0.412}.

\figitem{phase2/q24_cluster_choropleth.png}{1.0}{Q24 \textbullet\ Cluster choropleth.}{fig:q24a}
\figitem{phase2/q24_cluster_sizes.png}{0.95}{Q24 \textbullet\ Cluster population scale.}{fig:q24b}
\figitem{phase2/q24_radar.png}{0.95}{Q24 \textbullet\ Z-score signature of each cluster across the five features.}{fig:q24c}
\figitem{phase2/q24_parallel_coords.png}{1.0}{Q24 \textbullet\ Parallel coordinates showing each district's path through the feature space, coloured by cluster.}{fig:q24d}
\figitem{phase2/q24_priority_bar.png}{0.95}{Q24 \textbullet\ Cluster-level opportunity ranking.}{fig:q24e}
\figitem{phase2/q24_cagr_pop.png}{0.95}{Q24 \textbullet\ Cluster CAGR vs population.}{fig:q24f}

\textbf{Why this many visuals.} K-Means is a multidimensional method
and cannot be communicated with one chart. The choropleth shows
geography; the radar and parallel coordinates show feature signatures;
the priority bar converts the analysis into a decision artefact.

\textbf{Reading.} Four clusters:
\begin{itemize}
  \item \kw{Formal-Served Core} — 31 districts, CAGR 0.19\%, \stat{5.6M residents}.
  \item \kw{Peripheral Growth} — 19 districts, CAGR 2.83\%, \stat{9.2M residents}.
  \item \kw{Mixed} — 12 districts, CAGR 2.17\%, \stat{4.2M residents}.
  \item \kw{Low-Activity Outskirts} — 6 districts, CAGR 11.93\%.
\end{itemize}
The first three clusters together = \stat{62 districts, 18.9 million
residents}. \kw{This is the addressable market for any product that
bridges the two systems.}

% --------------------------------------------------------------------
\subsection{Q25 \textbullet\ Masari bridge — informal-to-formal connectability}

\textbf{Why.} The natural follow-up to Q24: \emph{can} an informal stop
be bridged to a formal node, and how far is the bridge? This is the
descriptive layer underneath the product story.

\textbf{Data.} Phase 1 boarding (informal-heavy stops) + integrated
formal layer (GTFS, metro, LRT, BRT).

\textbf{Method.} For each informal-heavy Phase 1 stop, find the
nearest formal node; compute distance; visualise the connector
network.

\figitem{phase2/q25_bridge_schematic.png}{1.0}{Q25 \textbullet\ The Masari bridge schematic. Each connector is informal-stop $\to$ nearest formal node. Short connectors are immediate routing wins; long connectors are where rider-contributed routing has to fill in.}{fig:q25}

\textbf{Reading.} A majority of informal-heavy stops sit within walking
distance of \emph{some} formal node — the bridges exist; they are just
not visible in any single application today. Section~\ref{sec:synthesis}
returns to this in the product story.
